\documentclass{../latex/utmreport}
\newcommand{\ReportTitle}{Lab 6: Inside the xv6 File System}
% Shared title-page and PDF metadata for every laboratory report.
\newcommand{\StudentName}{Gafenco Victor}
\newcommand{\StudentGroup}{FAF-241}
\newcommand{\CourseName}{Operating Systems}
\newcommand{\InstructorName}{Patricia Reițman}
\newcommand{\InstructorTitle}{Assistant lecturer}
\newcommand{\ReportCity}{Chișinău}
\newcommand{\ReportYear}{2026}
\newcommand{\CourseRepository}{https://github.com/yoda-dinmd/os-labs}

\selectlanguage{english}
\utmlogo{../latex/assets/UTM_Logo_English.jpeg}
\institution{Technical University of Moldova}
\faculty{Faculty of Computers, Informatics, and Microelectronics}
\department{Department of Software Engineering}
\studyprogramme{Software Engineering Study Programme}
\documenttype{Laboratory report}
\student{\StudentName}{\StudentGroup}
\supervisor{\InstructorName}{\InstructorTitle}
\renewcommand{\UTMCoordinatorLabel}{Instructor}
\cityyear{\ReportCity}{\ReportYear}
\reportsubtitle{\CourseName}
\hypersetup{pdfauthor={\StudentName},pdfsubject={\CourseName},
  pdftitle={\ReportTitle}}
\reporttitle{\ReportTitle}
\addto\captionsenglish{\renewcommand{\figurename}{Figure}\renewcommand{\tablename}{Table}\renewcommand{\bibname}{References}}
\renewcommand{\figurename}{Figure}
\renewcommand{\tablename}{Table}
\renewcommand{\bibname}{References}
\usepackage{float}
\usepackage{placeins}
\setlength{\emergencystretch}{2em}
\newcommand{\code}[1]{\texttt{\detokenize{#1}}}
% Keep screenshots beside their observations, with bounded height and width.
\newcommand{\screenshot}[4][\textwidth]{%
  \begin{figure}[H]
    \centering
    \includegraphics[width=#1,height=0.62\textheight,keepaspectratio]{screenshots/#2}
    \caption{#3}\label{#4}
  \end{figure}}
\newcommand{\reportreferences}[2]{%
  \clearpage\phantomsection\addcontentsline{toc}{chapter}{References}
  \begin{thebibliography}{9}
    \bibitem{handout} Technical University of Moldova, FCIM / FAF.
    \textit{#1}. Operating Systems laboratory handout, academic year 2026--2027.
    Supplied course document: \code{#2}.
  \end{thebibliography}}

\providecommand{\tightlist}{\setlength{\itemsep}{0pt}\setlength{\parskip}{0pt}}
\begin{document}
\makeutmreporttitlepage
\tableofcontents
\clearpage
\begin{utmabstract}[Abstract]
This laboratory examines the on-disk structures and visible behaviour of the xv6 filesystem. Two user programs expose inode metadata and recursively search directory trees. Fresh-image experiments demonstrate hard links, fixed-length names, directory entry sizes and inode identity. Source inspection explains block reclamation and storage limits, while traversal tests exercise descriptor management and the optional inode display.
\utmkeywords{xv6, inode, directory, hard link, recursive traversal, filesystem}
\end{utmabstract}
\utmintroduction

This laboratory examines directory entries, inode metadata, hard links
and storage limits in xv6. Two user programs were implemented:
\texttt{mystat}, which prints inode information, and \texttt{find},
which recursively searches by the last pathname component. The optional
\texttt{-n} form of \texttt{find} also prints inode numbers.

The work uses the same course revision as Lab 4,
{\small\texttt{06aad25c735fd3159bdfae5680be4eab7b1668b2}}, with the RISC-V GCC
toolchain and QEMU on openSUSE Tumbleweed. A separate working copy was
built from \texttt{\textasciitilde{}/xv6-riscv}. Both experiments began
with freshly built filesystem images: the first covers inspection and
link removal, and the second contains the exact directory tree
prescribed for \texttt{find}. QEMU was exited with Ctrl-A followed by X
between runs.

The source files are submitted as \texttt{user/mystat.c} and
\texttt{user/find.c}, together with the full Makefile and a focused
Makefile patch. The command transcripts and relevant kernel-source
excerpts are retained in \texttt{evidence/}.

\chapter{Reading the filesystem
listing}\label{reading-the-filesystem-listing}

\section{File types and root
identity}\label{file-types-and-root-identity}

The first \texttt{ls} produced the following selected entries:

\noindent\begin{minipage}{\linewidth}
\begin{lstlisting}
$ ls
.              1 1 1024
..             1 1 1024
README         2 2 2441
cat            2 7 35368
console        3 27 0
\end{lstlisting}
\end{minipage}\par\medskip

The columns are name, type, inode number and size in bytes.
\texttt{kernel/stat.h} defines \textbf{1 as a directory}
(\texttt{T\_DIR}), \textbf{2 as a regular file} (\texttt{T\_FILE}) and
\textbf{3 as a device} (\texttt{T\_DEVICE}). The console is therefore a
device inode rather than a regular file holding console contents.

The root directory has inode \textbf{1}, matching \texttt{ROOTINO} in
\texttt{kernel/fs.h}. At root, \texttt{.} denotes the current directory
and \texttt{..} denotes its parent, which is root itself. Both therefore
identify inode 1. The executable and console inode numbers depend on the
files included in the image; the handout's sample numbers are not fixed
requirements.

\section{An empty directory}\label{an-empty-directory}

\noindent\begin{minipage}{\linewidth}
\begin{lstlisting}
$ mkdir a
$ ls a
.              1 28 32
..             1 1 1024
$ mystat a
a: directory, inode 28, links 1, size 32 bytes
\end{lstlisting}
\end{minipage}\par\medskip

A directory entry consists of a two-byte \texttt{ushort\ inum} and a
14-byte name array, so \texttt{sizeof(struct\ dirent)} is \textbf{16
bytes}. Directory \texttt{a} contains two entries, \texttt{.} and
\texttt{..}, accounting for \texttt{2\ *\ 16\ =\ 32} bytes. Its
\texttt{.} entry identifies inode 28, and \texttt{..} identifies the
parent, inode 1. The 1024 displayed beside \texttt{..} is the parent
directory's size, not the length of that single entry.

\chapter{Implementing mystat and observing hard
links}\label{implementing-mystat-and-observing-hard-links}

\section{Interface and error
handling}\label{interface-and-error-handling}

\texttt{mystat} loops over every pathname operand and calls the
user-library \texttt{stat} function. Its type helper maps the three
defined types to \texttt{directory}, \texttt{file} and \texttt{device},
returning \texttt{?} for an unknown value. Successful calls print the
requested name, type, inode number, link count and size. Size is printed
with \texttt{\%lu}, supported by this revision's \texttt{user/printf.c},
preserving the \texttt{uint64} value in \texttt{struct\ stat}.

A failed lookup prints a diagnostic to descriptor 2 and processing
continues with later operands. The program exits with status 1 if any
lookup failed, and status 0 otherwise. Calling it without operands
prints a usage message and exits with status 1. This extends the
handout's basic loop with useful argument and failure reporting.

The Makefile retains the earlier user utilities and adds tab-indented
continuation lines for \texttt{\$U/\_mystat} and \texttt{\$U/\_find} in
\texttt{UPROGS}. A clean build of \texttt{kernel/kernel} and
\texttt{fs.img} completed successfully.

\section{Metadata before and after
linking}\label{metadata-before-and-after-linking}

\noindent\begin{minipage}{\linewidth}
\begin{lstlisting}
$ mystat README . console
README: file, inode 2, links 1, size 2441 bytes
.: directory, inode 1, links 2, size 1024 bytes
console: device, inode 27, links 1, size 0 bytes
$ ln README r2
$ mystat README r2
README: file, inode 2, links 2, size 2441 bytes
r2: file, inode 2, links 2, size 2441 bytes
$ rm README
$ mystat r2
r2: file, inode 2, links 1, size 2441 bytes
$ cat r2 | wc
48 336 2441
\end{lstlisting}
\end{minipage}\par\medskip

Creating \texttt{r2} adds a second directory entry pointing to inode 2
and increases its link count from 1 to 2. Removing \texttt{README}
decreases the count to 1. The remaining name still identifies inode 2,
the file remains 2441 bytes long, and its contents remain readable:
\texttt{wc} counted 48 lines, 336 words and 2441 bytes. A hard link does
not copy the data.

Root's link count is 2 in this run because \texttt{mkdir\ a} already
added a child directory. This differs from the handout's initial example
with link count 1. xv6 deliberately does not count a directory's own
\texttt{.} entry as a separate link; directory creation increases the
parent's count for the child's \texttt{..} entry. Therefore an empty xv6
directory has link count 1 even though it contains two directory
entries.

Error handling was also observed:

\noindent\begin{minipage}{\linewidth}
\begin{lstlisting}
$ mystat missing r2
mystat: cannot stat missing
r2: file, inode 2, links 1, size 2441 bytes
$ mystat
usage: mystat <path> [path ...]
\end{lstlisting}
\end{minipage}\par\medskip

\section{Unlinking and freeing
blocks}\label{unlinking-and-freeing-blocks}

In the pinned \texttt{kernel/sysfile.c}, the file's link count is
decremented by \textbf{\texttt{ip-\textgreater{}nlink-\/-;} at line 244}
in \texttt{sys\_unlink}. The following \texttt{iupdate(ip)} persists the
inode metadata. The earlier \texttt{dp-\textgreater{}nlink-\/-} at line
239 instead updates a parent when a directory is removed; these two
counters must not be confused.

The file's data blocks are freed only after its link count is zero and
its final in-memory reference is being released. In
\texttt{kernel/fs.c}, \texttt{iput} tests
\texttt{ip-\textgreater{}ref\ ==\ 1\ \&\&\ ip-\textgreater{}valid\ \&\&\ ip-\textgreater{}nlink\ ==\ 0}
at line 357, then calls \texttt{itrunc(ip)} at line 365. \texttt{itrunc}
frees direct data blocks, indirect data blocks and the indirect-address
block. This revision then calls \texttt{ifree} to make the inode
reusable. An open descriptor can delay this reclamation after the last
directory name has been removed. After only \texttt{rm\ README}, the
surviving \texttt{r2} link already prevents reclamation.

\section{What an on-disk inode
stores}\label{what-an-on-disk-inode-stores}

\texttt{struct\ dinode} in \texttt{kernel/fs.h} contains:

{
\begin{longtable}[]{@{}
  >{\raggedright\arraybackslash}p{0.47\linewidth}
  >{\raggedright\arraybackslash}p{0.47\linewidth}@{}}
\hline
\begin{minipage}[b]{\linewidth}\raggedright
Field
\end{minipage} & \begin{minipage}[b]{\linewidth}\raggedright
Meaning
\end{minipage} \\
\hline
\endhead
\hline
\endlastfoot
\texttt{short\ type} & Directory, regular file, device, or unused
inode \\
\texttt{short\ major} & Device major number, for a device inode \\
\texttt{short\ minor} & Device minor number, for a device inode \\
\texttt{short\ nlink} & Number of counted directory links \\
\texttt{uint\ size} & Logical size in bytes \\
\texttt{uint\ addrs{[}NDIRECT\ +\ 1{]}} & Twelve direct addresses and
one indirect-block address \\
\end{longtable}
}

The structure occupies 64 bytes: four two-byte fields, one four-byte
size and thirteen four-byte addresses. A filename lives in a directory
entry, not in this inode. This xv6 revision has no inode ownership, Unix
permission-mode fields, timestamps or symbolic-link type. The
\texttt{dev} and \texttt{ino} reported by \texttt{struct\ stat} identify
an inode's location; they are not fields stored inside
\texttt{struct\ dinode}.

\chapter{Measuring filesystem limits}\label{measuring-filesystem-limits}

\section{Fixed-size names}\label{fixed-size-names}

The constants in \texttt{kernel/fs.h} are \texttt{BSIZE\ =\ 1024},
\texttt{NDIRECT\ =\ 12}, \texttt{NINDIRECT\ =\ BSIZE\ /\ sizeof(uint)}
and \texttt{DIRSIZ\ =\ 14}.

\noindent\begin{minipage}{\linewidth}
\begin{lstlisting}
$ mkdir averyveryverylongname
$ ls
averyveryveryl 1 29 32
$ mystat averyveryverylongname averyveryveryl
averyveryverylongname: directory, inode 29, links 1, size 32 bytes
averyveryveryl: directory, inode 29, links 1, size 32 bytes
\end{lstlisting}
\end{minipage}\par\medskip

The \texttt{ls} excerpt selects the newly created directory. Its stored
name is \textbf{\texttt{averyveryveryl}}, the first 14 characters.
\texttt{DIRSIZ} explains the truncation. Path lookup likewise uses the
first 14 characters of a component, so both spellings identify inode 29.
A full-length 14-byte name need not contain a terminating NUL inside the
directory entry, which matters when passing it to string functions.

\section{Maximum file length}\label{maximum-file-length}

One indirect block holds \texttt{1024\ /\ 4\ =\ 256} block addresses.
Together with twelve direct blocks, the inode addresses 268 data blocks:

\noindent\begin{minipage}{\linewidth}
\begin{lstlisting}
NINDIRECT = 1024 / 4 = 256
MAXFILE   = 12 + 256 = 268 data blocks
maximum   = 268 * 1024 = 274432 bytes = 268 KiB
\end{lstlisting}
\end{minipage}\par\medskip

This is the structural maximum file-data length. The indirect-address
block consumes another disk block but does not add to the file's logical
contents. A nearly full disk may run out of free blocks before a file
reaches this limit. The calculation follows the source constants; it is
not a claim that a maximum-size file was written during this exercise.

\section{Disk layout from mkfs}\label{disk-layout-from-mkfs}

The clean build printed:

\noindent\begin{minipage}{\linewidth}
\begin{lstlisting}
nmeta 47 (boot, super, log blocks 31, inode blocks 13, bitmap blocks 1) blocks 1953 total 2000
balloc: first 1102 blocks have been allocated
\end{lstlisting}
\end{minipage}\par\medskip

The disk contains \textbf{2000 blocks}, each 1024 bytes, for
\textbf{2,048,000 bytes} overall. Its contiguous layout is:

{
\begin{table}[H]
\centering
\caption{Filesystem image layout}
\begin{tabular}{lll}
\hline
Region & Block numbers & Block count \\
\hline


Boot & 0 & 1 \\
Superblock & 1 & 1 \\
Log & 2--32 & 31 \\
Inodes & 33--45 & 13 \\
Free-block bitmap & 46 & 1 \\
Data area & 47--1999 & 1953 \\
\end{tabular}
\end{table}
}

The metadata total is \texttt{1\ +\ 1\ +\ 31\ +\ 13\ +\ 1\ =\ 47};
\texttt{2000\ -\ 47\ =\ 1953} blocks remain in the data region. The log
allocation includes its header. The 1953 data-area blocks are not all
initially free: the root directory, README, executables and any
indirect-address blocks already occupy part of them. Immediately after
this build, the first 1102 disk blocks were marked allocated, leaving
898 blocks unallocated before the boot-time and user-created changes.

\chapter{Implementing and testing
find}\label{implementing-and-testing-find}

\section{Traversal and resource
handling}\label{traversal-and-resource-handling}

\texttt{find\ \textless{}dir\textgreater{}\ \textless{}name\textgreater{}}
compares the final component of each visited pathname with the requested
name using \texttt{strcmp}. A matching directory is printed as well as
searched, and a matching regular file is printed without being
traversed. A starting file operand is also accepted. Trailing slashes on
the starting pathname are normalized except at root.

For a directory, the program reads fixed-size \texttt{struct\ dirent}
records, skips entries with \texttt{inum\ ==\ 0}, and excludes
\texttt{.} and \texttt{..} from recursion. Each name is copied into a
\texttt{DIRSIZ\ +\ 1} buffer and explicitly NUL-terminated. Child paths
include a separator only when required, and paths exceeding the 512-byte
buffer policy, including the terminator, produce a diagnostic.

Before recursion, the names are retained in a heap-allocated linked list
and the directory descriptor is closed. Child pathnames are also
allocated on the heap and freed after use. This avoids retaining one
open descriptor per ancestor: xv6's \texttt{NOFILE} limit is 16, with
standard descriptors already occupying three slots. Directory read
errors, failed lookup/open calls and allocation failures are reported to
standard error and result in a nonzero final status.

The traversal targets the stable directory trees in this experiment. Its
name snapshot is not an atomic filesystem snapshot; concurrent renames
can cause lookup errors. Recursion still consumes the finite xv6 user
stack, so the program does not claim unlimited depth.

\section{Required test on a fresh
image}\label{required-test-on-a-fresh-image}

After exiting QEMU, \texttt{make\ clean} and a new build created a fresh
image, removing the earlier README-link experiment. The required
commands then produced:

\noindent\begin{minipage}{\linewidth}
\begin{lstlisting}
$ mkdir a
$ mkdir a/b
$ echo x > a/b/r2
$ echo y > a/r2
$ ln cat r2
$ find . r2
./a/b/r2
./a/r2
./r2
$ find a b
a/b
\end{lstlisting}
\end{minipage}\par\medskip

All three matching files were found, and searching for \texttt{b} found
the matching directory. Excluding \texttt{.} and \texttt{..} prevents
cycling back to ancestors. The two nested files each contain a letter
and newline, while the root \texttt{r2} is a hard link to the executable
\texttt{cat}.

\section{Optional inode numbers}\label{optional-inode-numbers}

\noindent\begin{minipage}{\linewidth}
\begin{lstlisting}
$ find . r2 -n
./a/b/r2 30
./a/r2 31
./r2 7
$ mystat cat r2 a/r2 a/b/r2
cat: file, inode 7, links 2, size 35368 bytes
r2: file, inode 7, links 2, size 35368 bytes
a/r2: file, inode 31, links 1, size 2 bytes
a/b/r2: file, inode 30, links 1, size 2 bytes
\end{lstlisting}
\end{minipage}\par\medskip

Only \textbf{\texttt{./r2}} is the same file as \texttt{cat}: both
identify inode 7 and report two links. The other two paths have distinct
inodes, 30 and 31, even though all three last components are named
\texttt{r2}.

\section{Additional checks}\label{additional-checks}

A missing path produced \texttt{find:\ cannot\ stat\ missing}; invalid
argument counts and an unknown fourth argument produced the usage
message. A nonexistent target name produced no match lines. The
following cases also succeeded: \texttt{find\ a/\ r2},
\texttt{find\ a/r2\ r2}, \texttt{find\ /a\ r2}, a directory with a full
14-character name, and a directory containing a slot freed by removing a
file.

A deeper tree contained 19 nested directories, exceeding the per-process
descriptor count. Both ordinary and inode-number searches reached its
\texttt{needle} file without an open failure. A second tree contained 24
sibling directories; repeated searches found \texttt{wide/d23/needle}.
These results support the descriptor-closing strategy and the handling
of both deep and wide trees. Full commands and output are retained in
\texttt{evidence/5-edge-cases.txt} and
\texttt{evidence/6-depth-width.txt}.

\utmconclusions

The listing and \texttt{mystat} results connect names to inode
identities and link counts. Removing one hard link leaves the inode and
its contents reachable through another, while source inspection explains
when the final unlink can reclaim storage. The filesystem constants
account for 14-byte names, 16-byte directory entries and a 268 KiB
maximum file size. Finally, \texttt{find} turns the directory-entry
representation into a working recursive traversal, and its inode display
distinguishes equal names from references to the same file.
\end{document}
